\scanpage{187}
des $G_{\mathbf R}$-modules de dim finie sur $\mathbf R$. Le cône qui exprime cette effectivité est \uline{engendré} par les \deleted{$\unsure{u(s')|v(s')}$} rayons définis par les $v(s')$. \unsure{C'est} un fait qui exprime \inserted{sous forme ``analytique''} le fait que les \unsure{frobénius|frobenius|Frobenius} $f(s')\in G_{\mathbf R}^{\xi(s')}(\mathbf R)$ ont la vertu que leurs valeurs propres, dans \unsure{tte} représentation de $G_{\mathbf R}^{\xi(s')}\deleted{(\mathbf R)}$ -- et en fait $=$ de $G_{\deleted{\mathbf R}\mathbf C}^{\xi(s')}$ -- sont des \uline{entiers} algébriques, ce qui est une condition arithmétique assez draconienne !
\margin{186}
